\section{A path code with exponentially many ordered costs}
\label{sec:path-code}

Use $r$ path coordinates together with two distinguished logical
coordinates. Thus the logical group is
$G=\mathbb F_2^{r+2}$. Let $a_j$ be the following $\ell=2r+1$
columns in this group: the two distinguished unit vectors, all $r$
path unit vectors, and the $r-1$ sums of consecutive path unit vectors.
Their respective rail-size increments are
\[
 \delta_1=2Ar,\quad \delta_2=4Ar,\quad
 \delta_{\rm loop}=4,\quad\delta_{\rm edge}=2,
\]
where $A=2^{20}$ is fixed. Choose a sufficiently large fixed even
integer $K$, and set
\[
 R=KA^2r,\quad n_0=\ell R+1,\quad
 S=6r-2,\quad D_\delta=6Ar+S.
\]
Each short rail has size $R$ and each long rail size $R+\delta_j$.
The retained transversal indexed by $s\in G$ chooses its long rail
in pair $j$ exactly when $a_j\cdot s=1$. Including the shared carrier,
its odd degree is
\begin{equation}\label{eq:path-degree}
 n_s=n_0+2Ars_1+4Ars_2+4|s_{\rm path}|
                  +2|Ds_{\rm path}|,
\end{equation}
where $D$ is the adjacent-difference map on the path. In particular
$R\ge6\max_j\delta_j$, and all degrees are $\Theta(r^2)$.

\subsection{Positive higher-order coefficients and one dominant pair}

For odd integer arguments, the hardtanh cost has the analytic extension
\[
 w(z)=\sqrt{\pi z}\frac{\Gamma((z+1)/2)}{\Gamma(z/2)}
 =\sqrt{\frac\pi2}\left(z-\frac14+\frac1{32z}+O(|z|^{-2})\right).
\]
Put $c=\sqrt{\pi/2}/32$. The remainder is holomorphic and uniformly
$O(|z|^{-2})$ in a sufficiently far right half-plane. Cauchy's
estimate on a circle of radius $a(1-1/k)$ proves, for all
$2\le k\le\ell$ and $a\ge n_0$, once $K$ is large enough,
\begin{equation}\label{eq:path-derivative}
 \frac{ck!}{2a^{k+1}}\le(-1)^kw^{(k)}(a)
                  \le\frac{3ck!}{2a^{k+1}}.
\end{equation}
Indeed the relative remainder bound is $O(k^2/a)$, uniformly in
these indices, and $n_0\ge KA^2\ell^2/3$.

First regard the $\ell$ long-rail choices as independent bits $t_j$.
If $\widetilde W(T)$ is the Walsh coefficient of
$w(n_0+\sum_j\delta_jt_j)$, the finite-difference formula gives
\begin{equation}\label{eq:path-independent-bounds}
 \frac{ck!\prod_{j\in T}\delta_j}
      {2^{k+1}(n_0+D_\delta)^{k+1}}
 \le\widetilde W(T)\le
 \frac{3ck!\prod_{j\in T}\delta_j}{2^{k+1}n_0^{k+1}}
 \qquad(k=|T|\ge2).
\end{equation}
All singleton coefficients are negative because $w$ is increasing.
Restriction to the linear code pushes these coefficients forward:
\begin{equation}\label{eq:path-pushforward}
 W(v)=\sum_{T:\,\sum_{j\in T}a_j=v}\widetilde W(T).
\end{equation}
Only zero and the individual column directions can receive a
constant or singleton contribution. Every other nonzero $W(v)$ is
strictly positive.

Let $u=e_1+e_2$ in the distinguished coordinates. It is not a column,
and $W(u)\ge\widetilde W(\{1,2\})>0$. The total of the other
positive coefficients obeys
\begin{equation}\label{eq:path-dominance}
 \sum_{v\ne0,u}\max(W(v),0)\le\eta W(u),\qquad \eta<10^{-3},
\end{equation}
by increasing the fixed $K$ if necessary. To see this quantitatively,
the other independent degree-two products sum to at most
$6ArS+S^2/2$, whereas the distinguished product is $8A^2r^2$.
Their ratio in \eqref{eq:path-independent-bounds} is at most
$3(6ArS+S^2/2)(1+D_\delta/n_0)^3/(8A^2r^2)=O(A^{-1})$.
For $k\ge3$, use
$k!\sum_{|T|=k}\prod_{j\in T}\delta_j\le D_\delta^k$ and sum
the resulting geometric series. Relative to the distinguished pair
this contributes $O(D_\delta^3/(n_0\delta_1\delta_2))
=O((KAr)^{-1})$. Pushforward cannot increase the total positive mass
outside zero and $u$, so these bounds prove the claim.

\subsection{Many distinct effective costs}

For a path vector $v\in\mathbb F_2^r$, identify $(1,1,v)\in G$ and
write
\[
 \nu_v=W(1,1,v)-W(0,0,v).
\]
Averaging the distinguished bits gives the exact formula
\begin{equation}\label{eq:path-effective}
 \nu_v=-\frac12\left[\widehat w_{n_0+\delta_1}(v)
                         +\widehat w_{n_0+\delta_2}(v)\right],
\end{equation}
where $\widehat w_a(v)$ is the path Walsh coefficient of
$w(a+4|s|+2|Ds|)$. If $|v|\ge q\ge5$, every set of path columns
whose sum is $v$ has at least $h=\lceil q/2\rceil$ members.
Applying \eqref{eq:path-independent-bounds} just to the path columns
therefore proves
\begin{equation}\label{eq:path-tail}
 0<-\nu_v\le\frac{3c}{n_0}
                  \left(\frac{S}{2n_0}\right)^h.
\end{equation}
The geometric-series factor is absorbed using $S/(2n_0)<1/2$.
Strict positivity follows because the path loop columns span the path
space. This argument uses finite differences only; it does not assume
an unproved positive Laplace representation for $w$.

Define the signed degree-count polynomial
\[
 P_v(X)=\sum_{s\in\mathbb F_2^r}(-1)^{s\cdot v}
                              X^{2|s|+|Ds|}.
\]
At least $2^{r/2}$ of these polynomials are distinct. Here is an
explicit transfer-matrix certificate. For $e\in(0,1)$ set
$a=2e/(1+e^2)$ and
\[
 Z_v(a,e)=\sum_{y\in\mathbb F_2^{r-1}}e^{|y|}a^{|v+By|},
\]
where $B$ maps a set of path edges to its odd-degree vertex set.
The Fourier expansion of $\exp[-t(4|s|+2|Ds|)]$, with
$e=\tanh t$ and $a=\tanh2t$, is a common positive factor times
$Z_v(a,e)$. It is also $2^{-r}P_v(e^{-2t})$.
Let
\[
 M_0=\begin{pmatrix}1&ae\\a&e\end{pmatrix},\quad
 M_1=\begin{pmatrix}a&e\\1&ae\end{pmatrix},\quad
 c_0=\binom1a,\quad c_1=\binom a1.
\]
Directly summing the successive edge variables gives
\[
 \binom{Z_v}{Z_{v+e_1}}=M_{v_1}\cdots M_{v_{r-1}}c_{v_r}.
\]
At $e=1/2$, $a=4/5$, the first matrix sends a component ratio $x$ to
$f(x)=(x+2/5)/((4/5)x+1/2)$, and the second sends it to $1/f(x)$.
The interval $I=[4/5,5/4]$ is invariant under both maps, and
$f(I)=[20/19,11/10]\subset(1,\infty)$, whereas its reciprocal lies
below one. Both maps are injective. The resulting ratio therefore
recovers $v_1$, then $v_2$, and so forth, and the terminal ratios
$4/5$ and $5/4$ recover $v_r$. Thus all $2^r$ ordered pairs
$(Z_v,Z_{v+e_1})$ are distinct. If there were fewer than $2^{r/2}$
scalar values, there would be fewer than $2^r$ ordered pairs.
Distinct scalar values imply distinct $P_v$.

Distinct $P_v$ give distinct effective costs $\nu_v$ in
\eqref{eq:path-effective}. Indeed the two degree intervals, shifted
by $\delta_1$ and $\delta_2$, are disjoint because
$\delta_2-\delta_1=2Ar>S$. All degrees in both intervals lie in an
interval of length $D_\delta<2\sqrt{n_0}$ when $K$ is large.
Two distinct integers of the same squarefree part differ by more than
twice the square root of the smaller. Hence these distinct degrees
have distinct squarefree parts. Since $w(n)$ is a nonzero rational
multiple of $\sqrt n$ at every odd integer $n$, independence of
squarefree radicals forces every coefficient in a cost difference to
vanish. The disjoint upper branch then forces $P_v=P_{v'}$.

\subsection{An affine live set and quantitative tails}

Set $q=\lceil r/64\rceil$. For all sufficiently large $r$,
\[
 \sum_{j<q}\binom rj\le2^{r/8},\qquad \binom rq\le2^{r/8}.
\]
Start with one representative of each distinct polynomial, and discard
representatives of weight below $q$. Assign to each remaining vector
a $q$-subset of its support, and its two colors
\[
 \sum_{i=1}^r v_i\pmod2,\qquad
 \sum_{i\ \mathrm{odd}}v_i\pmod2.
\]
Pigeonholing over at most $4\binom rq$ choices gives an affine fiber
of size at least $2^{3r/8-3}$ in distinct-cost representatives. In
particular it contains at least $2^{r/3}$ for large $r$. Let $U$ be
this affine fiber, with its distinguished coordinates fixed to $(1,1)$.
All its path vectors have weight at least $q$. Put
$U'=U\cup(U+u)$. Every path loop or edge column changes at least
one of the two colors, and every distinguished loop changes exactly
one distinguished bit. Thus
\[
 U'\cap(U'+a_j)=\varnothing\quad\text{for every physical column }a_j.
\]
Also $0,u\notin U'$ and $u+a_j\notin U'$ for large $r$, since
vectors in $U'$ have at least $q>2$ path ones. Retain one representative
of each desired distinct effective cost as a live candidate, and
initially reveal the other fibers in $U'$ if necessary. The extra
initial observations only impose the additional nonpositive-sign
constraints already allowed in the rounding argument.

Finally \eqref{eq:path-tail} is $\exp[-\Omega(r\log r)]$, whereas
all geometric, norm-stability, and Fourier-rounding constants are
$\exp[O(r)]\operatorname{poly}(r)$. It therefore satisfies their
strict absorption inequalities for all sufficiently large $r$.
All parameters $A,K$ are fixed before taking this limit.

\subsection{Teacher resources}

The logical target is the sign of
$\sum_{i=1}^{r+2}\chi_{A_{e_i}}-(r+1)\chi_{A_0}$ and has margin one.
An odd parity of degree $n$ has an exact hardtanh spline representation
using at most $n+1$ neurons: transitions centered at the successive
midpoints of its attainable coordinate sums, and one constant neuron.
Incoming weights have magnitude one, biases at most $n$, and
readouts at most two. The teacher above therefore has width at most
\[
 C_r=(r+3)(n_0+1)+6Ar+8r-4=\Theta(r^3),
\]
since the sum of path basis-vector degree increments is $8r-4$.
Its readout weights are at most $2(r+1)$ and all coefficients are
bounded by $C_r$. The physical input dimension is
$d_r=1+2\ell R+D_\delta=\Theta(r^2)$.
The at least $2^{r/3}$ designated mistakes are consequently
$\exp(\Omega(C_r^{1/3}))$.
